\section{Spin systems and signed monomial duality}
\label{sec:monomial-duality}

\begin{discussion}
  \item Fix only the background needed later. Explain why the monomial
  basis is natural for local functions on \(\{0,1\}^{\Lambda}\).
  \item Distinguish assumptions used throughout from the polynomial
  growth and uniform pure-death assumptions introduced in
  \cref{sec:common-limit}.
\end{discussion}

\subsection{Spin systems and convergence terminology}

\begin{definition}[Spin system]
\label{def:spin-system}
Let \(\Lambda\) be countable and
\(\Omega=\{0,1\}^{\Lambda}\). For \(i\in\Lambda\), let \(\eta^i\)
denote the configuration obtained by flipping \(\eta(i)\). A spin
system is a continuous-time Markov process whose generator acts on
local functions by
\begin{equation}
  \cL f(\eta)
  =
  \sum_{i\in\Lambda}
  c_i(\eta)\bigl(f(\eta^i)-f(\eta)\bigr),
  \label{eq:spin-generator}
\end{equation}
where \(c_i(\eta)\ge0\). Its semigroup is denoted by
\((P_t)_{t\ge0}\).
\end{definition}

\begin{assumption}[Standing locality assumptions]
\label{ass:standing-locality}
The flip rates are uniformly bounded and finite-range. More
precisely, \(c_i\) depends on a finite set
\(N_*(i)=N(i)\cup\{i\}\), where \(i\notin N(i)\), and the resulting
graphical constructions started from a finite dual set are locally
finite.
\end{assumption}

\begin{definition}[Invariant measure]
\label{def:invariant-measure}
A probability measure \(\pi\) on \(\Omega\) is invariant for
\((P_t)_{t\ge0}\) if
\[
  \pi P_t=\pi
  \qquad(t\ge0).
\]
\end{definition}

\begin{definition}[Ergodicity]
\label{def:ergodicity}
The spin system is ergodic if there is a probability measure \(\pi\)
such that
\[
  \nu P_t\Longrightarrow\pi
  \qquad(t\to\infty)
\]
for every initial probability measure \(\nu\) on \(\Omega\). The limit
\(\pi\) is then the unique \invariantmeasure.
\end{definition}

\begin{definition}[Uniform exponential ergodicity on local functions]
\label{def:uniform-exponential-ergodicity}
The spin system is uniformly exponentially ergodic on local functions
if it has an \invariantmeasure\ \(\pi\) and there is \(\gamma>0\) such
that, for every local function \(f\), some \(C_f<\infty\) satisfies
\[
  \sup_{\eta\in\Omega}
  \abs{P_tf(\eta)-\pi(f)}
  \le C_fe^{-\gamma t}
  \qquad(t\ge0).
\]
\end{definition}

\subsection{Rate coefficients and the signed dual}

\begin{notation}[Monomials and rate coefficients]
\label{not:monomial-coefficients}
For \(A\Subset\Lambda\), set
\[
  \chi_A(\eta)=\prod_{i\in A}\eta(i),
  \qquad
  \chi_\vn=1.
\]
For \(x\in\{0,1\}\), put \(c_i^x(\eta)=c_i(\eta^{i,x})\).
The unique multilinear expansion on the neighbour set is
\begin{equation}
  c_i^x(\eta)
  =
  \sum_{S\subseteq N(i)}c_i^x(S)\chi_S(\eta).
  \label{eq:rate-expansion}
\end{equation}
Define
\begin{equation}
  a_i^\delta(S)=c_i^0(S),
  \qquad
  a_i^\beta(S)=-c_i^0(S)-c_i^1(S).
  \label{eq:signed-coefficients}
\end{equation}
\end{notation}

\begin{definition}[Signed additive set process]
\label{def:signed-dual}
Set
\[
  \delta_i(S)=\abs{a_i^\delta(S)},
  \qquad
  \beta_i(S)=\abs{a_i^\beta(S)}
  \quad(S\ne\vn),
  \qquad
  \beta_i(\vn)=0,
\]
and let
\(\sigma_i^\delta(S),\sigma_i^\beta(S)\in\{+,-\}\) be the
corresponding signs. For \(Y=(A,\sigma)\) and \(i\in A\), define
\[
  D_{i,S}Y
  =
  \bigl((A\setminus\{i\})\cup S,
        \sigma\sigma_i^\delta(S)\bigr)
\]
and, for \(S\ne\vn\),
\[
  B_{i,S}Y
  =
  \bigl(A\cup S,\sigma\sigma_i^\beta(S)\bigr).
\]
The generator of the signed additive set process is
\begin{align}
  \cD g(A,\sigma)
  ={}&
  \sum_{i\in A}\sum_{S\subseteq N(i)}
  \delta_i(S)
  \bigl(g(D_{i,S}(A,\sigma))-g(A,\sigma)\bigr)
  \notag\\
  &+
  \sum_{i\in A}\sum_{\vn\ne S\subseteq N(i)}
  \beta_i(S)
  \bigl(g(B_{i,S}(A,\sigma))-g(A,\sigma)\bigr).
  \label{eq:dual-generator}
\end{align}
Finally, set
\begin{equation}
  V(A)=\sum_{i\in A}V_i,
  \qquad
  V_i
  =
  \sum_{S\subseteq N(i)}\delta_i(S)
  +
  \sum_{\vn\ne S\subseteq N(i)}\beta_i(S)
  +
  a_i^\beta(\vn).
  \label{eq:fk-potential}
\end{equation}
\end{definition}

\begin{theorem}[Signed monomial duality]
\label{thm:monomial-duality}
Let \(H((A,\sigma),\eta)=\sigma\chi_A(\eta)\). Then
\[
  \cL_\eta H=\cD H+V(A)H.
\]
Consequently, if the dual starts from \((A,+)\), then
\begin{equation}
  P_t\chi_A(\eta)
  =
  \E_{(A,+)}\left[
    \sigma_t
    \exp\left(\int_0^tV(A_s)\,\dd s\right)
    \chi_{A_t}(\eta)
  \right].
  \label{eq:fk-duality}
\end{equation}
\end{theorem}

\begin{proof}
Only flips at sites \(i\in A\) affect \(\chi_A\). For such \(i\),
\[
  c_i(\eta)\bigl(\chi_A(\eta^i)-\chi_A(\eta)\bigr)
  =
  c_i^0(\eta)\chi_{A\setminus\{i\}}(\eta)
  -
  \bigl(c_i^0(\eta)+c_i^1(\eta)\bigr)\chi_A(\eta).
\]
Substitution of \eqref{eq:rate-expansion} identifies the off-diagonal
terms with \eqref{eq:dual-generator}; the difference between their
total rates and the diagonal term is \eqref{eq:fk-potential}. The
semigroup formula is the Feynman--Kac consequence of the generator
identity; see, for example, \citet{jansen-kurt2014}.
\end{proof}

\begin{construction}[Graphical construction]
\label{con:dual-graphical}
For every \(i\) and \(S\subseteq N(i)\), place independent Poisson
clocks of rate \(\delta_i(S)\). For every nonempty
\(S\subseteq N(i)\), also place independent clocks of rate
\(\beta_i(S)\). A clock interaction is applied precisely when its
source \(i\) is active. Adjoin the deterministic initial interaction
\[
  (\infty,0,\mathsf{init},A;\,+),
\]
whose target installs the initial active set and whose formal source
is not a lattice site.
\end{construction}
